% Paso 6: Hoja de Contexto Canónico
% Proyecto: Análisis Exegético Profundo
% Ministerio "La Gloria es del Señor"

\documentclass[11pt,a4paper]{article}
\usepackage[utf8]{inputenc}
\usepackage[spanish]{babel}
\usepackage[margin=2.5cm]{geometry}
\usepackage{graphicx}
\usepackage{fancyhdr}
\usepackage{xcolor}
\usepackage{tcolorbox}
\usepackage{enumitem}
\usepackage{tabularx}
\usepackage{multirow}
\usepackage{amsfonts}
\usepackage{fontspec}
\usepackage{titlesec}

% Configuración de colores
\definecolor{forestgreen}{RGB}{27,67,50}
\definecolor{sagegreen}{RGB}{45,90,39}
\definecolor{olivegreen}{RGB}{64,145,108}
\definecolor{mintgreen}{RGB}{82,183,136}
\definecolor{lightmint}{RGB}{116,198,157}
\definecolor{softcream}{RGB}{216,243,220}
\definecolor{papergray}{RGB}{118,128,150}

% Configuración de encabezados
\pagestyle{fancy}
\fancyhf{}
\fancyhead[L]{\textcolor{forestgreen}{\textbf{Análisis Exegético Profundo}}}
\fancyhead[R]{\textcolor{papergray}{Paso 6: Contexto Canónico}}
\fancyfoot[C]{\textcolor{papergray}{\thepage}}
\fancyfoot[R]{\textcolor{papergray}{\scriptsize Ministerio "La Gloria es del Señor"}}

% Configuración de títulos
\titleformat{\section}
{\color{forestgreen}\Large\bfseries}
{\thesection}{1em}{}

\titleformat{\subsection}
{\color{sagegreen}\large\bfseries}
{\thesubsection}{1em}{}

% Configuración de cajas
\tcbuselibrary{skins}
\newtcolorbox{infobox}[1]{
    colback=softcream,
    colframe=olivegreen,
    fonttitle=\bfseries,
    title=#1,
    left=10pt,
    right=10pt,
    top=10pt,
    bottom=10pt
}

\newtcolorbox{checklistbox}{
    colback=white,
    colframe=mintgreen,
    left=15pt,
    right=15pt,
    top=10pt,
    bottom=10pt,
    boxrule=2pt
}

\newtcolorbox{contextbox}{
    colback=lightmint!20,
    colframe=sagegreen,
    left=10pt,
    right=10pt,
    top=8pt,
    bottom=8pt,
    boxrule=1.5pt
}

\newtcolorbox{canonbox}{
    colback=olivegreen!10,
    colframe=olivegreen,
    left=10pt,
    right=10pt,
    top=8pt,
    bottom=8pt,
    boxrule=1.5pt
}

\begin{document}

% Título principal
\begin{center}
    {\Huge\color{forestgreen}\textbf{PASO 6}}\\[0.5cm]
    {\Large\color{sagegreen}\textbf{CONTEXTO CANÓNICO}}\\[0.3cm]
    {\large\color{papergray}El Texto dentro de toda la Escritura}\\[1cm]
    
    \begin{tcolorbox}[colback=olivegreen!10, colframe=olivegreen, width=0.8\textwidth]
        \centering
        \textcolor{forestgreen}{\textbf{Proyecto: Análisis Exegético Profundo}}\\
        \textcolor{papergray}{Metodología de 8 Pasos para Estudio Bíblico Riguroso}
    \end{tcolorbox}
\end{center}

\vspace{1cm}

% Información del estudiante
\begin{infobox}{Información del Proyecto}
\begin{tabularx}{\textwidth}{|X|X|}
\hline
\textbf{Nombre del Estudiante:} & \rule{6cm}{0.4pt} \\[0.8cm]
\hline
\textbf{Texto Bíblico:} & \rule{6cm}{0.4pt} \\[0.8cm]
\hline
\textbf{Libro Bíblico:} & \rule{3cm}{0.4pt} \quad \textbf{Capítulo:} \rule{2cm}{0.4pt} \\[0.8cm]
\hline
\textbf{Fecha:} & \rule{6cm}{0.4pt} \\[0.8cm]
\hline
\end{tabularx}
\end{infobox}

\section{Objetivo del Paso 6: Contexto Canónico}

El análisis del contexto canónico sitúa el pasaje dentro del marco de toda la revelación bíblica. Examinaremos el texto en relación con su contexto inmediato, el libro completo, y toda la Escritura para entender su lugar en el plan revelatorio de Dios.

\subsection{Descripción del Proceso}

Analizaremos el pasaje dentro de círculos concéntricos de contexto: párrafos vecinos, capítulo, libro completo, corpus del autor, testamento, y toda la Biblia. Buscaremos textos paralelos, temas relacionados y el desarrollo progresivo de la revelación.

\begin{infobox}{Principio de Interpretación}
\textbf{Scriptura Scripturam interpretatur} - "La Escritura interpreta a la Escritura"\\
El mejor intérprete de cualquier pasaje bíblico es la Biblia misma.
\end{infobox}

\section{Lista de Verificación}

\begin{checklistbox}
\textbf{Marca cada elemento completado:}

\begin{itemize}[leftmargin=1cm]
    \item[$\square$] Analizar contexto inmediato (párrafos vecinos)
    \item[$\square$] Considerar contexto del capítulo completo
    \item[$\square$] Examinar el lugar dentro del libro completo
    \item[$\square$] Buscar paralelos en otros libros del mismo autor
    \item[$\square$] Identificar temas teológicos relacionados
    \item[$\square$] Examinar desarrollo progresivo de la revelación
    \item[$\square$] Buscar citas o alusiones del AT/NT
    \item[$\square$] Considerar el testimonio unificado de la Escritura
\end{itemize}
\end{checklistbox}

\newpage

\section{Contexto Inmediato}

\subsection{Párrafos Circundantes}

\begin{contextbox}
\textbf{Párrafo Anterior:} \rule{4cm}{0.4pt} \textbf{Párrafo Posterior:} \rule{4cm}{0.4pt}
\end{contextbox}

\textbf{¿Cómo se conecta tu pasaje con el material que lo precede?}
\vspace{2cm}

\textbf{¿Cómo se conecta con el material que lo sigue?}
\vspace{2cm}

\textbf{¿Qué palabras clave, temas o ideas se conectan con el contexto inmediato?}
\vspace{2cm}

\subsection{Función en el Argumento Local}

\textbf{¿Qué función cumple tu pasaje en el argumento o narrativa del capítulo?}
\begin{itemize}[leftmargin=1cm]
    \item[$\square$] Introduce un tema nuevo
    \item[$\square$] Desarrolla un tema existente
    \item[$\square$] Proporciona evidencia o apoyo
    \item[$\square$] Ilustra con ejemplo
    \item[$\square$] Resume o concluye
    \item[$\square$] Hace transición
    \item[$\square$] Otro: \rule{6cm}{0.4pt}
\end{itemize}

\section{Contexto del Libro}

\subsection{Propósito del Libro}

\begin{canonbox}
\textbf{Propósito/Tema Principal del Libro:}\\[0.5cm]
\rule{12cm}{0.4pt}\\[0.5cm]
\textbf{Audiencia Original:} \rule{8cm}{0.4pt}\\[0.5cm]
\textbf{Ocasión de Escritura:} \rule{8cm}{0.4pt}
\end{canonbox}

\textbf{¿Cómo contribuye tu pasaje al propósito general del libro?}
\vspace{3cm}

\subsection{Estructura del Libro}

\textbf{¿En qué sección del libro está ubicado tu pasaje?}

\begin{tabularx}{\textwidth}{|c|X|X|}
\hline
\textbf{Sección} & \textbf{Capítulos} & \textbf{Tema/Propósito} \\
\hline
I & & \\[1cm]
\hline
II & & \\[1cm]
\hline
III & & \\[1cm]
\hline
IV & & \\[1cm]
\hline
V & & \\[1cm]
\hline
\end{tabularx}

\textbf{¿Tu pasaje está en:}
\begin{itemize}[leftmargin=1cm]
    \item[$\square$] Introducción/Prólogo
    \item[$\square$] Desarrollo principal
    \item[$\square$] Clímax del libro
    \item[$\square$] Conclusión/Epílogo
    \item[$\square$] Transición entre secciones
\end{itemize}

\newpage

\section{Paralelos Textuales}

\subsection{Paralelos en el Mismo Autor}

\textbf{¿Qué otros textos del mismo autor tratan temas similares?}

\begin{tabularx}{\textwidth}{|X|X|X|}
\hline
\textbf{Referencia} & \textbf{Tema/Concepto Paralelo} & \textbf{Similaridades/Diferencias} \\
\hline
& & \\[1.5cm]
\hline
& & \\[1.5cm]
\hline
& & \\[1.5cm]
\hline
& & \\[1.5cm]
\hline
\end{tabularx}

\subsection{Paralelos en Otros Libros}

\textbf{¿Qué textos de otros autores bíblicos desarrollan temas similares?}

\begin{tabularx}{\textwidth}{|X|X|X|}
\hline
\textbf{Referencia} & \textbf{Autor/Libro} & \textbf{Relación con tu Texto} \\
\hline
& & \\[1.5cm]
\hline
& & \\[1.5cm]
\hline
& & \\[1.5cm]
\hline
& & \\[1.5cm]
\hline
\end{tabularx}

\section{Citas y Alusiones}

\subsection{Referencias al Antiguo Testamento}

\begin{contextbox}
\textbf{Completa solo si tu texto es del Nuevo Testamento:}
\end{contextbox}

\textbf{¿Tu texto cita o alude a pasajes del AT?}

\begin{tabularx}{\textwidth}{|X|X|X|}
\hline
\textbf{Texto AT} & \textbf{Tipo de Uso} & \textbf{Propósito de la Referencia} \\
\hline
& Cita directa / Alusión & \\[1cm]
\hline
& Cita directa / Alusión & \\[1cm]
\hline
& Cita directa / Alusión & \\[1cm]
\hline
\end{tabularx}

\textbf{¿Cómo usa el autor NT el contexto original del AT?}
\vspace{2cm}

\subsection{Ecos Intertextuales}

\textbf{¿Hay ecos más sutiles de otros textos bíblicos?}
\vspace{2cm}

\textbf{¿Tu texto hace eco de algún patrón tipológico bíblico?}
\vspace{2cm}

\newpage

\section{Desarrollo Temático}

\subsection{Temas Teológicos Principales}

\textbf{¿Qué temas teológicos desarrolla tu texto?}

\begin{tabularx}{\textwidth}{|X|X|X|}
\hline
\textbf{Tema Teológico} & \textbf{Desarrollo en tu Texto} & \textbf{Desarrollo en Otra Escritura} \\
\hline
& & \\[1.5cm]
\hline
& & \\[1.5cm]
\hline
& & \\[1.5cm]
\hline
& & \\[1.5cm]
\hline
\end{tabularx}

\subsection{Progresión de la Revelación}

\textbf{¿Cómo se relaciona tu texto con el desarrollo progresivo de estos temas en la Escritura?}

\begin{itemize}[leftmargin=1cm]
    \item[$\square$] Introduce un tema nuevo
    \item[$\square$] Desarrolla revelación previa
    \item[$\square$] Clarifica o expande conceptos
    \item[$\square$] Provee cumplimiento de promesas/profecías
    \item[$\square$] Aplica verdades establecidas
    \item[$\square$] Otro: \rule{6cm}{0.4pt}
\end{itemize}

\textbf{Explica la progresión:}
\vspace{3cm}

\section{Unidad Canónica}

\subsection{Armonía con Toda la Escritura}

\textbf{¿Cómo armoniza tu interpretación con el testimonio total de la Escritura?}
\vspace{3cm}

\textbf{¿Hay alguna tensión aparente con otros textos que necesite resolverse?}
\vspace{3cm}

\textbf{¿Tu interpretación está en línea con doctrinas bíblicas claramente establecidas?}
\vspace{2cm}

\subsection{Contribución al Canon}

\textbf{¿Qué contribución única hace tu texto al canon bíblico?}
\vspace{3cm}

\textbf{¿Qué se perdería del mensaje bíblico sin este texto?}
\vspace{3cm}

\newpage

\section{Tipología y Cumplimiento}

\subsection{Elementos Tipológicos}

\textbf{¿Hay elementos tipológicos en tu texto?}
\begin{itemize}[leftmargin=1cm]
    \item[$\square$] Personas (tipos de Cristo, etc.)
    \item[$\square$] Eventos (éxodo, exilio, etc.)
    \item[$\square$] Instituciones (sacerdocio, templo, etc.)
    \item[$\square$] Objetos (serpiente de bronce, etc.)
    \item[$\square$] No aplica
\end{itemize}

\textbf{Si aplica, describe la relación tipológica:}
\vspace{3cm}

\subsection{Cumplimiento Profético}

\textbf{¿Tu texto cumple profecías previas o contiene profecías?}
\vspace{2cm}

\textbf{¿Hay cumplimiento inmediato y/o escatológico?}
\vspace{2cm}

\section{Síntesis Canónica}

\subsection{Lugar en la Historia de la Salvación}

\textbf{¿Dónde encaja tu texto en la historia general de la salvación?}
\vspace{3cm}

\textbf{¿Cómo contribuye al plan redentor de Dios?}
\vspace{3cm}

\subsection{Hallazgos del Contexto Canónico}

\textbf{¿Qué aspecto del análisis canónico más enriquece tu comprensión del texto?}
\vspace{3cm}

\textbf{¿Cambió algo en tu interpretación a la luz del contexto canónico?}
\vspace{3cm}

\begin{infobox}{Principio Hermenéutico}
La Biblia es tanto una biblioteca de 66 libros como un libro unificado. Cada texto debe entenderse tanto en su contexto particular como dentro del testimonio total de la Escritura.
\end{infobox}

\section{Siguiente Paso}

Con el análisis del contexto canónico completado, procederás al \textbf{Paso 7: Síntesis Interpretativa}, donde integraremos todos los hallazgos en una interpretación coherente.

\vspace{1cm}

\begin{center}
\begin{tcolorbox}[colback=mintgreen!20, colframe=mintgreen, width=0.7\textwidth]
\centering
\textcolor{forestgreen}{\textbf{¡Contexto Bíblico Establecido!}}\\
\textcolor{papergray}{Has situado tu texto dentro del gran panorama de la revelación bíblica.}
\end{tcolorbox}
\end{center}

\end{document}